\chapter{Datamodels}
\label{toc:datamodels}

\section{Syntax}
\label{toc:syntax}

A datamodel is defined with the clause block \texttt{datamodel}. It can
contain \texttt{pmodel}, which are ghost model fields common to all
constructors, \texttt{pinvariant}, which are global invariants for the
datamodel, and none or several constructors, defined with \texttt{pcase}.

A constructor Cons, defined with a clause block \texttt{pcase} must contain
a clause \texttt{pwhen} --- if the constructor is the last constructor defined
in the datamodel, the clause \texttt{pwhen} may be omitted --- that declares the
guard of the constructor. It can also
contain \texttt{pframe}, which are distinct local frames or self-framing local
locations, and \texttt{pinvariant}.

In function contracts, the instances of data models consumed and produced
are specified using the \texttt{produces} and \texttt{consumes} clauses.
Each global model field must be associated with a corresponding variable
or value. The instances of data models manipulated in a loop are specified
similarly using the \texttt{loop frame} clause.

The heap and the frame can be manipulated inside a function with several
clauses. The clause \texttt{heap} allows redefining the heap and the clause
\texttt{frame} allows renaming some instances of data models in the frame.
The clause \texttt{produce} is used to create a new instance of a data model in
the frame with a given constructor ; all the model fields must be associated
with a value and all the local frames must be associated with an instance of
a data model of the frame.
Conversely, the clause \texttt{consume} is used to remove an instance of a
data model of the frame with a given constructor. Similarly, to
\texttt{produce}, all the model fields and the local frame must be associated
with a value and a new instance of a data model.

\begin{figure}[htbp]
  \begin{center}
    \begin{tabular}{|ccl@{\hspace{.5cm}}l|}
    % LTeX: enabled=false
    \hline
    \rule{0pt}{2em}
    \textit{top-level-clause} & ::= & \ldots &\\
    & $\mid$ &
    \texttt{datamodel} \textit{id} \{&\\
                              &&\quad\textit{data-decl} &\\
                              &&\qquad\ldots   &\\
                              &&\quad\textit{data-decl} &\\
                              && \} & \\

    \rule{0pt}{2em}
      \textit{contract-clause} & ::= & \ldots &\\
      & $\mid$ &
      \texttt{consumes} \textit{term}\,, \ldots \;; & \\
      & $\mid$ &
      \texttt{produces} \textit{term}\,, \ldots \;; & \\
      \rule{0pt}{2em}
      \textit{loop-clause} & ::= & \ldots &\\
      & $\mid$ &
      \texttt{loop frame} \textit{term}\,, \ldots \;; & \\

    \rule{0pt}{2em}
      \textit{code-annotation-clause} & ::= & \ldots &\\
    & $\mid$ &
      \texttt{heap} \textit{term}\,, \dots \;;&\\
    & $\mid$ &
      \texttt{frame} \textit{term}\,, \dots \;;&\\
    & $\mid$ &
      \texttt{call} \textit{assignation} ;&\\
    & $\mid$ &
      \texttt{consume} \textit{id}: \textit{assignation} ; &\\
    & $\mid$ &
      \texttt{produce} \textit{id}: \textit{assignation} ; &\\

    \rule{0pt}{2em}
      \textit{assignation} & ::= &
      \{ \textit{id} \textbackslash{with}
         \textit{.id} = \textit{term} , \dots \;\} &\\

    \rule{0pt}{2em}
      \textit{data-decl} & ::= &
      \texttt{pmodel} (\textit{logic-type}) \textit{id}, \ldots \;; &\\
    & $\mid$ &
      \texttt{pframe} \textit{term}, \ldots \;; &\\
    & $\mid$ &
      \texttt{pinvariant} \textit{term}, \dots \;;&\\
    & $\mid$ &
      \texttt{pcase} \textit{id} \{&\\
       &&\quad\textit{case-decl} &\\
       &&\qquad\ldots            &\\
       &&\quad\textit{case-decl} &\\
       &&\}
    &\\

    \rule{0pt}{2em}
      \textit{case-decl} & ::= &
    \texttt{pwhen} \textit{term} ; &\\
    & $\mid$ &
    \texttt{pframe} \textit{term}, \ldots \;; &\\
    & $\mid$ &
    \texttt{pinvariant} \textit{term}, \dots \;; &\\

    &&& \\
    \hline
  % LTeX: enabled=true
  \end{tabular}
  \end{center}
  \caption{Syntax of Datamodels}
  \label{fig:datamodels}
\end{figure}
